\section*{Hoofdstuk \ref{ch:b1:radiation-heat-transport} --- Straling en warmtetransport}
% values: calc/ch27.py

\begin{solution}{exo:b1:radiation-heat-transport:1}
(a)~$S_0/\pi = 1361/\pi = \qty{433.2}{W.m^{-2}}$. (b)~Aan de evenaar bij een equinox
$\bar Q = (S_0/\pi)(a/r)^2$, dus $(a/r)^2 = 436.7/433.2 = 1.0080$ en $r/a =
1.0080^{-1/2} = 0.9960$: de aarde staat \qty{0.4}{\percent} dichter bij de zon dan
haar gemiddelde afstand. (c)~$\bar Q = 433.2 \times 0.9683 \times \cos\ang{23.44} =
\qty{384.9}{W.m^{-2}}$.
\end{solution}

\begin{solution}{exo:b1:radiation-heat-transport:2}
(a)~$415.6/172.4 = 2.41$. (b)~De zone tussen de parallellen $\pm\varphi$ heeft de
oppervlakte $2\pi a_\oplus^2 \times 2\sin\varphi$, een fractie $\sin\varphi$ van de bol:
$\sin\ang{37.0} = 0.602$, \qty{60}{\percent}. (c)~Op een as in \omterm{def:b1:spherical-astronomy:geographic}{breedtegraad} lijken de
banden nabij de polen even breed als die nabij de evenaar, maar zij bestrijken veel
minder oppervlak: de zone poolwaarts van \ang{60} is slechts $1 - \sin\ang{60} =
\qty{13.4}{\percent}$ van het oppervlak. Tegen $\sin\varphi$ vullen de tropen het
grootste deel van de as.
\end{solution}

\begin{solution}{exo:b1:radiation-heat-transport:3}
(a)~$T_e = (\text{OLR}/\sigma)^{1/4}$: \qty{256.4}{K}, \qty{241.3}{K} en
\qty{254.5}{K}. (b)~Geabsorbeerd zonlicht: $287.7/75.8 = 3.80$; OLR: $244.9/192.3 =
1.27$. (c)~De \omterm{def:b1:radiation-heat-transport:net-radiation}{netto straling} volgt het geabsorbeerde zonlicht: positief en groot in de
tropen, sterk negatief op hoge breedtegraden. De OLR, die weinig varieert, compenseert
de daling van het geabsorbeerde zonlicht niet.
\end{solution}

\begin{solution}{exo:b1:radiation-heat-transport:4}
(a)~$2.82 \times 365.25 = \qty{1030}{mm}$, \qty{1.03}{m} per jaar. (b)~$1.030 \times
78.1 = \qty{80.4}{W.m^{-2}}$.
(c)~De $LE$ van de reanalyse is \qty{2}{W.m^{-2}} groter. Over vele jaren kan de atmosfeer
geen waterdamp blijven opstapelen: al het verdampte water komt weer naar beneden, dus moeten $LE$ en $LP$ in het wereldgemiddelde gelijk zijn. Het verschil is een onbalans van het model, net als zijn stralingsonbalans.
\end{solution}

\begin{solution}{exo:b1:radiation-heat-transport:5}
(a)~$\bar Q = S_0(a/r)^2\sin\delta = 1361 \times 0.9683 \times 0.3978 =
\qty{524.2}{W.m^{-2}}$. (b)~$1361 \times 1.0334 \times 0.3978 =
\qty{559.5}{W.m^{-2}}$; de verhouding is $1.0334/0.9683 = 1.067$. (c)~De pool heeft
zonlicht 24 uur van de 24 op een constante hoogte, dus is haar gemiddelde $\sin h_\odot$ over
de dag $\sin\delta = 0.398$. Aan de evenaar staat de zon om het middaguur hoog maar is zij de helft van de tijd onder; het gemiddelde van $\sin h_\odot$ over de 24 uur is $\cos\delta/\pi =
0.292$. De verhouding is $\pi\tan\delta = 1.362$
(\cref{prop:b1:radiation-heat-transport:polar-summer}).
\end{solution}

\begin{solution}{exo:b1:radiation-heat-transport:6}
(a)~$\cos h_0 = -\tan\ang{40}\tan\ang{23.44} = -0.3638$, $h_0 = \ang{111.33} =
\qty{1.943}{rad}$; de dag duurt $24 \times 111.33/180 = \qty{14.84}{h}$; $\bar Q =
(1361/\pi) \times 0.9683 \times (1.943 \times 0.6428 \times 0.3978 + 0.7660 \times
0.9175 \times 0.9315) = \qty{483.0}{W.m^{-2}}$. (b)~$\cos h_0 = +0.3638$, $h_0 =
\ang{68.67}$, \qty{9.16}{h}; $\bar Q = \qty{155.9}{W.m^{-2}}$. (c)~De verhouding is
$483.0/155.9 = 3.10$. De afstand draagt $0.9683/1.0334 = 0.937$ bij: zij werkt tegen de noordelijke zomer in. De meetkunde, de hoogte van de zon en de lengte van de dag, draagt $3.10/0.937 = 3.31$ bij, het hele seizoenscontrast en meer.
\end{solution}

\begin{solution}{exo:b1:radiation-heat-transport:7}
(a)~$B = H/LE$: \num{0.093} aan de evenaar, \num{0.374} op \ang{25} N,
\num{0.338} op \ang{45} N, \num{-0.152} op \ang{70} N. (b)~\ang{25} N: het zonale
gemiddelde bevat de grote hete woestijnen, droog land met weinig water om te verdampen, waar de
verwarmde bodem de lucht in plaats daarvan opwarmt. (c)~$H < 0$ betekent dat warmte stroomt van de
lucht naar het oppervlak: in het jaargemiddelde zijn de sneeuw en het zee-ijs kouder dan de lucht
erboven, die wordt opgewarmd door lucht die van lagere breedtegraden is aangevoerd.
\end{solution}

\begin{solution}{exo:b1:radiation-heat-transport:8}
(a)~Beoordeeld: $\alpha = 100/340 = 0.294$, $R = 340 - 100 - 239 =
\qty{+1}{W.m^{-2}}$ (de \num{0.7} afgerond). Reanalyse: $\alpha = 117.0/341.9 =
0.342$, $R = 341.9 - 117.0 - 237.7 = \qty{-12.8}{W.m^{-2}}$ met deze afgeronde
termen, \qty{-12.9}{W.m^{-2}} met de onafgeronde van
\cref{ex:b1:radiation-heat-transport:global}. (b)~De reanalyse
absorbeert $224.8$ tegen $240$, \qty{15.2}{W.m^{-2}} minder, omdat zij \qty{17}{W.m^{-2}}
meer terugkaatst; haar OLR is slechts \qty{1.3}{W.m^{-2}} lager. De extra terugkaatsing verklaart de onbalans. (c)~$R$ van de Zuidpool tot de
Noordpool integreren geeft de wereldintegraal, $-12.9 \times 4\pi a_\oplus^2 = -12.9 \times
5.10\times10^{14} = \qty{-6.6}{PW}$. Een transport over de Noordpool is
onmogelijk: er ligt niets voorbij, de parallel is een punt. De propositie
geldt alleen als $\int R\,\mathrm dA = 0$, wat de stationaire toestand vereist.
\end{solution}

\begin{solution}{exo:b1:radiation-heat-transport:9}
(a)~Het oppervlak is $2\pi a_\oplus^2 \times 2\sin\ang{30} = 2\pi a_\oplus^2 =
\qty{2.55e14}{m^2}$, de helft van de aarde; het overschot is $33 \times 2.55\times10^{14} =
\qty{8.4}{PW}$. (b)~\qty{4.2}{PW} over elke parallel. (c)~De twee tellen op tot
\qty{8.5}{PW}, het overschot, maar zij zijn ongelijk: het zuidelijke halfrond exporteert
\qty{0.5}{PW} meer dan het noordelijke. Het overschot wordt niet gedeeld op de evenaar: in deze reanalyse steekt \qty{0.36}{PW} van het overschot van de noordelijke tropen de evenaar zuidwaarts over.
\end{solution}

\begin{solution}{exo:b1:radiation-heat-transport:10}
(a)~Schrijf het dagelijks gemiddelde als $\bar Q = (a/r)^2 F(\varphi, \delta)$, met $F$ de
haakjes van \cref{thm:b1:radiation-heat-transport:daily} maal $S_0/\pi$. Gedurende
$\mathrm dt$ is de ontvangen energie per vierkante meter $\bar Q\,\mathrm dt$, en
$\mathrm dt = r^2\,\mathrm d\lambda/h$, dus $\bar Q\,\mathrm dt = (a^2/h)
F(\varphi, \delta(\lambda))\,\mathrm d\lambda$: de factor $r^2$ valt weg.
(b)~Over een jaar is de energie $(a^2/h)\int_0^{2\pi} F(\varphi,
\delta(\lambda))\,\mathrm d\lambda$. De formule blijft gelijk wanneer $\varphi$ en
$\delta$ samen van teken veranderen, $F(-\varphi, \delta) = F(\varphi, -\delta)$, en
$\delta(\lambda + \pi) = -\delta(\lambda)$, omdat $\sin\delta = \sin\varepsilon
\sin\lambda$. Het verschuiven van $\lambda$ met $\pi$ in de integraal toont dat $-\varphi$
dezelfde energie ontvangt als $\varphi$. (c)~De noordelijke lente en zomer duren
\num{7.6} dagen langer, omdat de aarde dan nabij het aphelium staat en langzaam beweegt; de
zuidelijke zomer is helderder maar korter. Volgens (a) heffen de twee effecten elkaar precies op.
\end{solution}

\begin{solution}{exo:b1:radiation-heat-transport:11}
(a)~$\langle (a/r)^2\rangle = \frac1P\int_0^P \frac{a^2}{r^2}\,\mathrm dt =
\frac1P\int_0^{2\pi}\frac{a^2}{r^2}\frac{r^2}{h}\,\mathrm d\lambda = \frac{2\pi
a^2}{hP} = \frac{2\pi a^2}{2\pi a b} = \frac ab$, en $b = a\sqrt{1-e^2}$.
(b)~De momentane gemiddelde waarde over de wereld $\frac14 S_0(a/r)^2$ middelt uit tot $\frac14 S_0
\langle(a/r)^2\rangle = S_0/(4\sqrt{1-e^2})$. (c)~$340.25/\sqrt{1 - 0.0167^2} =
\qty{340.30}{W.m^{-2}}$, \qty{0.05}{W.m^{-2}} of \qty{0.014}{\percent} boven
$S_0/4$.
\end{solution}

\begin{solution}{exo:b1:radiation-heat-transport:12}
(a)~Atmosfeer $0.78 \times 4.10 = \qty{3.20}{PW}$, oceaan $0.22 \times 4.10 =
\qty{0.90}{PW}$. (b)~De atmosfeer alleen blijkt meer te dragen, \qty{5.0}{PW}, dan de atmosfeer en de oceaan van de reanalyse samen, \qty{4.02}{PW}: het totale transport van de reanalyse is tekort met minstens een vijfde ($1 - 4.02/5.0 = 0.20$).
Haar \omterm{def:b1:radiation-heat-transport:net-radiation}{netto straling} varieert te weinig tussen de tropen en de polen. (c)~Het verwijderen van een constante uit $R$ maakt haar wereldintegraal nul, maar laat haar variatie met de \omterm{def:b1:spherical-astronomy:geographic}{breedtegraad} ongewijzigd; de fout zit in die variatie, in de albedo van de wolken en oppervlakken van het model op elke \omterm{def:b1:spherical-astronomy:geographic}{breedtegraad}, niet alleen in het wereldgemiddelde.
\end{solution}

\begin{solution}{pb:b1:radiation-heat-transport:1}
\begin{enumerate}
  \item Juni: $\cos h_0 = -0.3638$, $h_0 = \ang{111.33}$, $\bar Q =
        \qty{483.0}{W.m^{-2}}$. December: $h_0 = \ang{68.67}$, $\bar Q =
        \qty{155.9}{W.m^{-2}}$. Maart: $h_0 = \ang{90}$, $\bar Q = (1361/\pi)
        \times 1.0081 \times \cos\ang{40} = \qty{334.6}{W.m^{-2}}$.
  \item Juni: $\tan\ang{80}\tan\ang{23.44} = 2.46 > 1$, poolse dag, $h_0 = \pi$:
        $\bar Q = 1361 \times 0.9683 \times \sin\ang{80}\sin\ang{23.44} =
        \qty{516.3}{W.m^{-2}}$. December: poolnacht, $\bar Q = 0$. Maart:
        $(1361/\pi) \times 1.0081 \times \cos\ang{80} = \qty{75.8}{W.m^{-2}}$.
  \item Op \ang{80} N gaat de zon niet onder ($h_0 = \ang{180}$, 24 uur), terwijl zij op \ang{40} N \qty{14.8}{h} schijnt. Haar hoogte is lager, tussen
        \ang{13.4} en \ang{33.4}, maar de term $h_0\sin\varphi\sin\delta$, die groeit met de lengte van de dag en de \omterm{def:b1:spherical-astronomy:geographic}{breedtegraad}, wint.
  \item De bandwaarden liggen \qty{0.5}{\percent}, \qty{0.8}{\percent}, \qty{1.7}{\percent}, \qty{1.7}{\percent} en
        \qty{1.3}{\percent} boven de centrale waarden. Ten eerste gebruikt de reanalyse een iets grotere zonne-instraling (vraag~5), \qty{0.5}{\percent} meer. Ten tweede is een bandgemiddelde niet de waarde in het midden: de evenaarwaartse helft van elke band heeft meer oppervlak en meer zonlicht dan de poolwaartse helft.
  \item Volgens \cref{prop:b1:radiation-heat-transport:global}, $S_0 = 4\sqrt{1-e^2}
        \times 341.9 = \qty{1367}{W.m^{-2}}$, \qty{0.47}{\percent} meer dan \qty{1361}{W.m^{-2}}.
  \item $A = 2\pi a_\oplus^2(1 - \sin\ang{40}) = 2\pi \times (6.371\times10^6)^2
        \times 0.3572 = \qty{9.110e13}{m^2}$, een fractie $(1 - \sin\ang{40})/2 =
        0.179$ van de aarde.
  \item $(1 - \alpha)\bar Q$: \num{198.0}, \num{159.3}, \num{124.7}, \num{92.5},
        \num{74.7} \unit{W.m^{-2}}.
  \item $R$: \num{-28.1}, \num{-51.0}, \num{-76.6}, \num{-104.9}, \num{-116.9}
        \unit{W.m^{-2}}. Alle vijf banden zijn in tekort; de poolband
        \ang{80}--\ang{90} heeft het grootste.
  \item $T_e = (\text{OLR}/\sigma)^{1/4}$: \num{251.3}, \num{246.8}, \num{244.1},
        \num{242.9}, \num{241.1} \unit{K}.
  \item Geabsorbeerd: $198.0/74.7 = 2.65$; OLR: $226.1/191.6 = 1.18$.
  \item De zon staat lager, en licht dat onder een scherpe hoek aankomt, wordt sterker teruggekaatst door wolken en water; sneeuw en zee-ijs, heldere oppervlakken, bedekken naar de pool toe meer van de bodem en de zee.
  \item De OLR gaat als $T_e^4$, en $T_e$ varieert over de kap met slechts \qty{10}{K}, \qty{4}{\percent}: de hoge breedtegraden worden veel warmer gehouden dan hun zonlicht alleen zou toelaten, door de warmte die uit het zuiden wordt aangevoerd.
  \item In \unit{m^2}: \num{3.143e13}, \num{2.550e13}, \num{1.879e13},
        \num{1.151e13}, \num{0.387e13}; de som is \qty{9.110e13}{m^2}.
  \item $R_iA_i$: \num{-0.883}, \num{-1.300}, \num{-1.439}, \num{-1.207},
        \num{-0.453} \unit{PW}; de kap: \qty{-5.28}{PW}.
  \item $\alpha_\text{kap} = \sum\alpha_i\bar Q_iA_i/\sum\bar Q_iA_i = 0.400$,
        tegen \num{0.342} voor de aarde.
  \item $-12.90 \times 5.10\times10^{14} = \qty{-6.58}{PW}$: de aarde zou warmte verliezen tegen die snelheid en afkoelen. De beoordeelde onbalans is
        $+0.7 \times 5.10\times10^{14} = \qty{+0.36}{PW}$, achttien keer kleiner en van het tegengestelde teken: de onbalans van de reanalyse is een modelfout.
  \item Gecorrigeerd $R_i + 12.90$: \num{-15.2}, \num{-38.1}, \num{-63.7},
        \num{-92.0}, \num{-104.0} \unit{W.m^{-2}}; vermogens (in PW) \num{-0.478},
        \num{-0.972}, \num{-1.197}, \num{-1.059}, \num{-0.403} \unit{PW}; de kap:
        \qty{-4.108}{PW}, een gemiddeld tekort van \qty{45.1}{W.m^{-2}}.
  \item De energie van de kap verandert niet van jaar tot jaar, dus moet wat zij aan de top van de atmosfeer verliest, erin komen, en de enige weg erin is over de 40e parallel: $T(\ang{40}) = -2\pi a_\oplus^2\int_{\ang{40}}^{\ang{90}}
        R\cos\varphi\,\mathrm d\varphi$.
  \item $T(\ang{40}\,\text{N}) = +\qty{4.11}{PW}$, noordwaarts.
  \item $LP = 2.4655\times10^6 \times 1.95/86400 = \qty{55.6}{W.m^{-2}}$, meer dan
        $LE = \qty{42.4}{W.m^{-2}}$: er valt meer water op de kap dan er van verdampt; het verschil wordt als damp aangevoerd.
  \item $(LE - LP) - 2.07 = 42.4 - 55.6 - 2.07 = \qty{-15.3}{W.m^{-2}}$; maal de
        oppervlakte van de kap: \qty{-1.40}{PW}. De atmosfeer brengt \qty{1.40}{PW}
        latente warmte over \ang{40} N.
  \item $1.40/4.11 = 0.34$: een derde.
  \item Oceaan $0.22 \times 4.11 = \qty{0.90}{PW}$; atmosfeer $0.78 \times 4.11 =
        \qty{3.20}{PW}$, waarvan \qty{1.40}{PW} latent en \qty{1.80}{PW}
        gedragen als voelbare warmte en potentiële energie van de lucht.
  \item Een ondergrens. Het totaal van de reanalyse op \ang{43} N, \qty{4.02}{PW}, is minder dan de atmosfeer alleen in de studie, \qty{5.0}{PW}; het ware totaal is groter.
  \item Ongeveer \qty{4.11}{PW}, \qty{4.11e15}{W}, volgens de straling van de reanalyse, een derde daarvan als latente warmte; de ware waarde is nog groter.
\end{enumerate}
\end{solution}
